% Options for packages loaded elsewhere
\PassOptionsToPackage{unicode}{hyperref}
\PassOptionsToPackage{hyphens}{url}
\documentclass[
]{article}
\usepackage{xcolor}
\usepackage{amsmath,amssymb}
\setcounter{secnumdepth}{-\maxdimen} % remove section numbering
\usepackage{iftex}
\ifPDFTeX
  \usepackage[T1]{fontenc}
  \usepackage[utf8]{inputenc}
  \usepackage{textcomp} % provide euro and other symbols
\else % if luatex or xetex
  \usepackage{unicode-math} % this also loads fontspec
  \defaultfontfeatures{Scale=MatchLowercase}
  \defaultfontfeatures[\rmfamily]{Ligatures=TeX,Scale=1}
\fi
\usepackage{lmodern}
\ifPDFTeX\else
  % xetex/luatex font selection
\fi
% Use upquote if available, for straight quotes in verbatim environments
\IfFileExists{upquote.sty}{\usepackage{upquote}}{}
\IfFileExists{microtype.sty}{% use microtype if available
  \usepackage[]{microtype}
  \UseMicrotypeSet[protrusion]{basicmath} % disable protrusion for tt fonts
}{}
\makeatletter
\@ifundefined{KOMAClassName}{% if non-KOMA class
  \IfFileExists{parskip.sty}{%
    \usepackage{parskip}
  }{% else
    \setlength{\parindent}{0pt}
    \setlength{\parskip}{6pt plus 2pt minus 1pt}}
}{% if KOMA class
  \KOMAoptions{parskip=half}}
\makeatother
\usepackage{longtable,booktabs,array}
\usepackage{calc} % for calculating minipage widths
% Correct order of tables after \paragraph or \subparagraph
\usepackage{etoolbox}
\makeatletter
\patchcmd\longtable{\par}{\if@noskipsec\mbox{}\fi\par}{}{}
\makeatother
% Allow footnotes in longtable head/foot
\IfFileExists{footnotehyper.sty}{\usepackage{footnotehyper}}{\usepackage{footnote}}
\makesavenoteenv{longtable}
\usepackage{graphicx}
\makeatletter
\newsavebox\pandoc@box
\newcommand*\pandocbounded[1]{% scales image to fit in text height/width
  \sbox\pandoc@box{#1}%
  \Gscale@div\@tempa{\textheight}{\dimexpr\ht\pandoc@box+\dp\pandoc@box\relax}%
  \Gscale@div\@tempb{\linewidth}{\wd\pandoc@box}%
  \ifdim\@tempb\p@<\@tempa\p@\let\@tempa\@tempb\fi% select the smaller of both
  \ifdim\@tempa\p@<\p@\scalebox{\@tempa}{\usebox\pandoc@box}%
  \else\usebox{\pandoc@box}%
  \fi%
}
% Set default figure placement to htbp
\def\fps@figure{htbp}
\makeatother
\setlength{\emergencystretch}{3em} % prevent overfull lines
\providecommand{\tightlist}{%
  \setlength{\itemsep}{0pt}\setlength{\parskip}{0pt}}
\usepackage{bookmark}
\IfFileExists{xurl.sty}{\usepackage{xurl}}{} % add URL line breaks if available
\urlstyle{same}
\hypersetup{
  hidelinks,
  pdfcreator={LaTeX via pandoc}}

\author{}
\date{}

\begin{document}

\section{论证过程：为什么选择加权/平均误差约束（基于原文）}\label{ux8bbaux8bc1ux8fc7ux7a0bux4e3aux4ec0ux4e48ux9009ux62e9ux52a0ux6743ux5e73ux5747ux8befux5deeux7ea6ux675fux57faux4e8eux539fux6587uxff09}

\subsection{1.
系统工作性质分析（引用题目原文）}\label{1-ux7cfbux7edfux5de5ux4f5cux6027ux8d28ux5206ux6790ux5f15ux7528ux9898ux76eeux539fux6587uxff09}

\begin{quote}
``某公司 AI 实验室长期运行一套中小型 AI
算力（用于所在城市噪声信号分析，基于傅里叶变换进行信号处理，搭配 8 块
GPU 组成计算集群）。''
\end{quote}

由此可以得到系统特点：

\begin{itemize}
\item
  \textbf{批量数据处理}：数据按小时/天累积分析，而非实时反馈控制
\item
  \textbf{非实时任务}：噪声分析结果不用于即时安全或控制决策
\item
  \textbf{任务量大且均匀}：每天固定总任务量，可动态调节 GPU
  负载和传输速率
\end{itemize}

\begin{quote}
结论：系统属于典型的
\textbf{批处理型数据分析任务}，对单小时或瞬时误差敏感性低，更关注整体误差或平均精度。
\end{quote}

\begin{center}\rule{0.5\linewidth}{0.5pt}\end{center}

\subsection{2.
对比误差约束方式}\label{2-ux5bf9ux6bd4ux8befux5deeux7ea6ux675fux65b9ux5f0f}

\begin{longtable}[]{@{}llll@{}}
\toprule\noalign{}
约束类型 & 定义 & 对调度的限制 & 是否符合任务性质 \\
\midrule\noalign{}
\endhead
\bottomrule\noalign{}
\endlastfoot
\textbf{瞬时误差约束} & 每小时误差 ≤ 5\% & 每小时 GPU
负载和传输速率必须落在严格可行区 &
不完全符合：浪费调度灵活性，功耗和电费优化受限 \\
\textbf{加权/平均误差约束} & 全天平均误差 ≤ 5\% &
可以在部分小时允许略高误差，只要平均满足 &
符合批处理任务：分析结果整体可靠，调度灵活，可优化功耗和电费 \\
\end{longtable}

\begin{quote}
核心点：噪声信号分析任务对
\textbf{整体统计结果敏感}，而非对每小时瞬时误差敏感。
\end{quote}

\begin{center}\rule{0.5\linewidth}{0.5pt}\end{center}

\subsection{3.
数据特性支撑（结合题目）}\label{3-ux6570ux636eux7279ux6027ux652fux6491ux7ed3ux5408ux9898ux76eeuxff09}

\begin{itemize}
\item
  噪声信号分析属于

  统计分析 + 信号处理

  ：

  \begin{itemize}
  \item
    原始数据存在随机波动
  \item
    单小时误差可能大，但整体频谱分析和日累计结果稳定
  \end{itemize}
\item
  文献支撑：
\end{itemize}

\begin{enumerate}
\def\labelenumi{\arabic{enumi}.}
\item
  Botchkarev, A. (2018)

  \begin{itemize}
  \item
    指出批量数据分析任务常采用 \textbf{平均误差或整体指标}
    评估精度，而不是严格瞬时误差
  \item
    支持采用加权/平均误差约束建模
  \end{itemize}
\item
  Jangam \& Pedda Muntala (2023)

  \begin{itemize}
  \item
    讨论分布式批处理系统误差管理，强调允许部分时间段误差略高，只要平均误差满足要求
  \item
    与噪声信号分析任务的性质相符
  \end{itemize}
\end{enumerate}

\begin{quote}
说明：采用加权/平均误差约束既符合题目数据特性，也符合行业惯例。
\end{quote}

\begin{center}\rule{0.5\linewidth}{0.5pt}\end{center}

\subsection{4.
建模与优化优势}\label{4-ux5efaux6a21ux4e0eux4f18ux5316ux4f18ux52bf}

\begin{itemize}
\item
  调度灵活性增加

  ：

  \begin{itemize}
  \item
    高峰电价时段可降低 GPU 负载或传输速率
  \item
    谷时段提高负载 → 平均误差仍满足 5\%
  \end{itemize}
\item
  功耗和电费优化空间更大

  ：

  \begin{itemize}
  \item
    避免每小时固定高负载 → 降低能源消耗
  \item
    可以结合电池储能策略，优化全天调度
  \end{itemize}
\item
  建模更贴合题目参考数据

  ：

  \begin{itemize}
  \item
    题目基准：GPU 60\%、传输 800 Mbps → 误差 30\%
  \item
    瞬时误差 ≤5\%几乎无法使用题目负载范围
  \item
    平均误差约束允许在题目范围内实现可行解
  \end{itemize}
\end{itemize}

\begin{center}\rule{0.5\linewidth}{0.5pt}\end{center}

\subsection{5. 综合结论}\label{5-ux7efcux5408ux7ed3ux8bba}

\begin{itemize}
\item
  \textbf{噪声信号分析任务性质}：误差约束应关注整体统计结果，而非每小时严格控制
\item
  采用加权/平均误差约束

  ：

  \begin{itemize}
  \item
    更符合批处理特性
  \item
    支持动态调度优化功耗和电费
  \item
    与行业实践一致（Botchkarev 2018; Jangam \& Pedda Muntala 2023）
  \end{itemize}
\item
  瞬时误差约束

  ：

  \begin{itemize}
  \item
    对任务调度过于严格
  \item
    可能导致调度策略不可行或浪费资源
  \end{itemize}
\end{itemize}

\begin{quote}
\textbf{建模决策}：在本题中，应采用 \textbf{全天平均误差 ≤5\%}
作为有效约束，并在报告中明确这是建模假设。
\end{quote}

\subsection{第二问：分时电价下的动态调度建模}\label{ux7b2cux4e8cux95eeux5206ux65f6ux7535ux4ef7ux4e0bux7684ux52a8ux6001ux8c03ux5ea6ux5efaux6a21}

\begin{quote}
核心逻辑说明：

如果采用瞬时误差约束，每小时都会被迫接近 \$G=85\$,
\$R=1200\$，动态调度基本消失。\\
因此采用处理量加权平均误差作为约束，更适合分时电价下的动态调度优化。
\end{quote}

\subsubsection{动态调度策略示意}\label{ux52a8ux6001ux8c03ux5ea6ux7b56ux7565ux793aux610f}

\subparagraph{峰时段：}\label{ux5cf0ux65f6ux6bb5}

可以降低 G 和 R\\
处理量 q\_t 变小\\
瞬时误差 E\_t 可能变大\\
但由于 q\_t 权重小，对全天平均误差影响有限

\subparagraph{谷时段：}\label{ux8c37ux65f6ux6bb5}

提高 G 和 R\\
处理量 q\_t 变大\\
瞬时误差 E\_t 降低\\
承担主要任务并拉低加权平均误差

\subparagraph{平时段：}\label{ux5e73ux65f6ux6bb5}

作为过渡和补充

\subsubsection{理论说明（论文可引用）}\label{ux7406ux8bbaux8bf4ux660eux8bbaux6587ux53efux5f15ux7528uxff09}

题目要求控制一天任务的整体分析误差。\\
由于不同小时完成的处理量不同，误差对最终结果的影响也应按该小时处理量加权。\\
因此本文采用处理量加权平均误差作为精度约束，而不是对每小时瞬时误差逐点约束。

\subsubsection{约束体系建议}\label{ux7ea6ux675fux4f53ux7cfbux5efaux8bae}

\begin{aligned}
60 \le G_t &\le 100, \\
800 \le R_t &\le 1200, \\
\sum_t q_t &\ge 1, \\
\frac{\sum_t q_t \cdot E_t}{\sum_t q_t} &\le 5
\end{aligned}

\subsubsection{目标函数}\label{ux76eeux6807ux51fdux6570}

\[\min \sum_t \text{price}_t \cdot P_t\]

其中：

\[P_t = P_\text{gpu}(G_t) + P_\text{trans}(R_t) + P_\text{cool}(G_t)\]

\subsubsection{计算结果}\label{ux8ba1ux7b97ux7ed3ux679c}

根据上述模型求解，第二问得到的最优动态调度策略如下：

\begin{longtable}[]{@{}llllllll@{}}
\toprule\noalign{}
时段类型 & 电价（元/kWh） & GPU负载 & 数据传输速率 & 系统功率 & 小时电费
& 每小时处理量 & 瞬时误差率 \\
\midrule\noalign{}
\endhead
\bottomrule\noalign{}
\endlastfoot
谷时段 & 0.8 & 100.000\% & 1200 Mbps & 8.6500 kW & 6.9200 元 & 0.062500
& 2.0000\% \\
平时段 & 1.2 & 78.642\% & 1200 Mbps & 5.7345 kW & 6.8814 元 & 0.049151 &
6.2716\% \\
峰时段 & 2.0 & 60.000\% & 1200 Mbps & 3.0500 kW & 6.1000 元 & 0.037500 &
10.0000\% \\
\end{longtable}

全天综合指标如下：

\begin{longtable}[]{@{}ll@{}}
\toprule\noalign{}
指标 & 数值 \\
\midrule\noalign{}
\endhead
\bottomrule\noalign{}
\endlastfoot
最小日电费 & 162.3362 元 \\
问题一固定方案日电费 & 190.3680 元 \\
电费降低 & 28.0318 元 \\
电费降低率 & 14.73\% \\
一天总能耗 & 150.2135 kWh \\
总处理量 & 1.2398 \\
加权平均误差 & 5.0000\% \\
\end{longtable}

将三类时段展开为 24 小时调度表：

\begin{longtable}[]{@{}lllllll@{}}
\toprule\noalign{}
小时 & 时段类型 & 电价 & GPU负载 & 数据传输速率 & 系统功率 & 小时电费 \\
\midrule\noalign{}
\endhead
\bottomrule\noalign{}
\endlastfoot
1 & 谷时段 & 0.8 & 100.000\% & 1200 Mbps & 8.6500 kW & 6.9200 元 \\
2 & 谷时段 & 0.8 & 100.000\% & 1200 Mbps & 8.6500 kW & 6.9200 元 \\
3 & 谷时段 & 0.8 & 100.000\% & 1200 Mbps & 8.6500 kW & 6.9200 元 \\
4 & 谷时段 & 0.8 & 100.000\% & 1200 Mbps & 8.6500 kW & 6.9200 元 \\
5 & 谷时段 & 0.8 & 100.000\% & 1200 Mbps & 8.6500 kW & 6.9200 元 \\
6 & 谷时段 & 0.8 & 100.000\% & 1200 Mbps & 8.6500 kW & 6.9200 元 \\
7 & 平时段 & 1.2 & 78.642\% & 1200 Mbps & 5.7345 kW & 6.8814 元 \\
8 & 平时段 & 1.2 & 78.642\% & 1200 Mbps & 5.7345 kW & 6.8814 元 \\
9 & 平时段 & 1.2 & 78.642\% & 1200 Mbps & 5.7345 kW & 6.8814 元 \\
10 & 峰时段 & 2.0 & 60.000\% & 1200 Mbps & 3.0500 kW & 6.1000 元 \\
11 & 峰时段 & 2.0 & 60.000\% & 1200 Mbps & 3.0500 kW & 6.1000 元 \\
12 & 平时段 & 1.2 & 78.642\% & 1200 Mbps & 5.7345 kW & 6.8814 元 \\
13 & 平时段 & 1.2 & 78.642\% & 1200 Mbps & 5.7345 kW & 6.8814 元 \\
14 & 平时段 & 1.2 & 78.642\% & 1200 Mbps & 5.7345 kW & 6.8814 元 \\
15 & 平时段 & 1.2 & 78.642\% & 1200 Mbps & 5.7345 kW & 6.8814 元 \\
16 & 平时段 & 1.2 & 78.642\% & 1200 Mbps & 5.7345 kW & 6.8814 元 \\
17 & 平时段 & 1.2 & 78.642\% & 1200 Mbps & 5.7345 kW & 6.8814 元 \\
18 & 平时段 & 1.2 & 78.642\% & 1200 Mbps & 5.7345 kW & 6.8814 元 \\
19 & 峰时段 & 2.0 & 60.000\% & 1200 Mbps & 3.0500 kW & 6.1000 元 \\
20 & 峰时段 & 2.0 & 60.000\% & 1200 Mbps & 3.0500 kW & 6.1000 元 \\
21 & 平时段 & 1.2 & 78.642\% & 1200 Mbps & 5.7345 kW & 6.8814 元 \\
22 & 平时段 & 1.2 & 78.642\% & 1200 Mbps & 5.7345 kW & 6.8814 元 \\
23 & 谷时段 & 0.8 & 100.000\% & 1200 Mbps & 8.6500 kW & 6.9200 元 \\
24 & 谷时段 & 0.8 & 100.000\% & 1200 Mbps & 8.6500 kW & 6.9200 元 \\
\end{longtable}

\subsubsection{计算结果分析}\label{ux8ba1ux7b97ux7ed3ux679cux5206ux6790}

从结果可以看出，最优方案具有明显的分时调度特征。谷时段电价最低，模型将
GPU 负载提高到 100\%，使该时段承担更多处理量；峰时段电价最高，模型将 GPU
负载降低到下限 60\%，减少高电价下的系统功率；平时段负载为
78.642\%，起到补充处理量和调节平均误差的作用。

该结果说明，第二问中的动态调度不是简单地追求每小时功率最低，而是把功率消耗重新分配到更便宜的电价时段。谷时段单小时功率最高，为
8.6500 kW，但由于电价只有 0.8 元/kWh，小时电费为 6.9200
元；峰时段单小时功率最低，为 3.0500 kW，即使电价达到 2.0
元/kWh，小时电费也只有 6.1000
元。由此可见，动态调度的核心是让``高功率''尽量出现在``低电价''时段。

峰时段瞬时误差为 10.0000\%，高于
5\%，但这并不违反模型约束。因为问题二采用的是处理量加权平均误差，而不是逐小时瞬时误差约束。峰时段每小时处理量只有
0.037500，对全天误差的权重较小；谷时段每小时处理量为
0.062500，且瞬时误差只有
2.0000\%，承担了更多低误差处理任务。最终全天加权平均误差为：

\[\bar{E}=
\frac{\sum_t q_tE_t}{\sum_t q_t}
=5.0000\%\]

这说明误差约束刚好取到边界，是限制最优解继续降低电费的关键约束。

与问题一固定方案相比，问题二的日电费从 190.3680 元降低到 162.3362
元，节约 28.0318 元，节约率为 14.73\%。同时，问题二总处理量为
1.2398，仍然大于
1，满足全天任务量要求。因此，问题二方案在任务量和误差约束均满足的条件下，实现了电费的明显下降。

\subsubsection{数据传输速率始终取上限的原因}\label{ux6570ux636eux4f20ux8f93ux901fux7387ux59cbux7ec8ux53d6ux4e0aux9650ux7684ux539fux56e0}

计算结果中，峰、平、谷三个时段的数据传输速率均为 1200
Mbps。这一现象不是模型没有动态调度，而是由题目给定的误差模型和传输功耗模型共同决定。

误差模型为：

\[E_t = 30 - 0.2(G_t - 60) + 0.05(800 - R_t)\]

当数据传输速率从 800 Mbps 提高到 1200 Mbps 时，误差变化为：

\[0.05(800-1200)=-20\]

即误差率下降 20 个百分点。而传输功率模型为：

\[P_{\text{trans}}(R)=0.25+0.0004(R-800)\]

因此：

\[P_{\text{trans}}(800)=0.25\text{ kW}, \quad
P_{\text{trans}}(1200)=0.41\text{ kW}\]

传输速率提高到上限只增加 0.16 kW 功率，却能显著降低误差。相比之下，提高
GPU 负载不仅会增加 GPU
本身功耗，还会增加冷却功耗。因此，在本题参数下，提高传输速率是一种``低功耗代价、高误差收益''的调节方式。模型会优先将数据传输速率推到上限，再通过
GPU 负载实现峰、平、谷时段之间的动态调度。

\subsubsection{图像结果分析}\label{ux56feux50cfux7ed3ux679cux5206ux6790}

以下为五张问题二图像的文字分析。论文排版时可将对应图片插入到相应段落之后。

\paragraph{图1：24
小时电价与调度曲线}\label{ux56fe124-ux5c0fux65f6ux7535ux4ef7ux4e0eux8c03ux5ea6ux66f2ux7ebf}

\pandocbounded{\includegraphics[keepaspectratio,alt={}]{C:/Users/Nuwanda/AppData/Roaming/Typora/typora-user-images/image-20260520000648673.png}}

该图展示 24 小时内 GPU
负载和数据传输速率的变化，并用背景色区分峰、平、谷时段。

图中 GPU 负载呈明显的阶梯变化：谷时段为 100\%，平时段约为
78.642\%，峰时段为
60\%。这与分时电价完全对应，说明模型确实将更多计算任务安排到低电价时段，而在高电价时段降低负载。数据传输速率曲线始终保持在
1200
Mbps，说明传输速率主要承担降低误差的作用，而不是作为电价响应的主要调度变量。

该图可以用于说明问题二的核心调度逻辑：GPU
负载负责响应电价变化，数据传输速率负责保持整体误差水平。

\paragraph{图2：24
小时功率与小时电费图}\label{ux56fe224-ux5c0fux65f6ux529fux7387ux4e0eux5c0fux65f6ux7535ux8d39ux56fe}

\pandocbounded{\includegraphics[keepaspectratio,alt={}]{C:/Users/Nuwanda/AppData/Roaming/Typora/typora-user-images/image-20260520000703996.png}}

该图展示系统功率和每小时电费之间的关系。

从图中可以看出，谷时段系统功率最高，但由于电价最低，小时电费并没有显著超过其他时段；峰时段系统功率最低，抵消了高电价带来的成本压力。平时段功率和电费均处于中间水平。

这说明问题二优化的对象是总电费，而不是总功率。若只看功率，谷时段似乎消耗更多；但结合电价后，谷时段高功率反而是合理的，因为同样的计算任务在低价时段完成成本更低。

\paragraph{图3：问题一固定方案与问题二动态方案对比图}\label{ux56fe3ux95eeux9898ux4e00ux56faux5b9aux65b9ux6848ux4e0eux95eeux9898ux4e8cux52a8ux6001ux65b9ux6848ux5bf9ux6bd4ux56fe}

\pandocbounded{\includegraphics[keepaspectratio,alt={}]{C:/Users/Nuwanda/AppData/Roaming/Typora/typora-user-images/image-20260520000726191.png}}

该图比较问题一固定方案和问题二动态方案在日电费、日能耗、总处理量和加权误差四个指标上的差异。

对比结果表明，问题二动态方案显著降低了日电费：由 190.3680 元下降到
162.3362 元，降幅为
14.73\%。同时，动态方案仍然满足总处理量和加权平均误差约束。该图说明，引入分时电价后，优化目标从``功耗最低''转为``电费最低''，动态方案能够利用电价结构获得更低运行成本。

需要注意的是，问题二总处理量为
1.2398，略低于问题一固定方案的处理量，但仍然高于任务要求
1。这说明问题二并不是盲目追求更高处理量，而是在满足任务需求的前提下减少高电价时段的无效冗余。

\paragraph{图4：传输速率取上限原因图}\label{ux56fe4ux4f20ux8f93ux901fux7387ux53d6ux4e0aux9650ux539fux56e0ux56fe}

\pandocbounded{\includegraphics[keepaspectratio,alt={}]{C:/Users/Nuwanda/AppData/Roaming/Typora/typora-user-images/image-20260520000732863.png}}

该图用于解释为什么数据传输速率在所有时段都取 1200 Mbps。

图中传输功率随传输速率增加的幅度较小，而误差率随传输速率提高下降明显。也就是说，增加传输速率带来的功耗成本较低，但带来的误差改善很大。与提高
GPU 负载相比，提高传输速率是一种更经济的降误差方式。

因此，最优解中传输速率保持最大值并不矛盾。问题二的动态性主要体现在 GPU
负载随电价变化，而传输速率则由于较高的误差收益和较低的功耗代价，稳定保持在上限。

\subsubsection{第二问结论}\label{ux7b2cux4e8cux95eeux7ed3ux8bba}

综合计算结果和图像分析，第二问的最优调度策略可以概括为：

\begin{enumerate}
\def\labelenumi{\arabic{enumi}.}
\item
  谷时段提高 GPU 负载，承担更多处理量；
\item
  峰时段降低 GPU 负载，减少高电价时段电费；
\item
  平时段保持中等负载，平衡任务量和误差；
\item
  数据传输速率始终保持 1200 Mbps，以较小功耗代价降低整体误差；
\item
  全天加权平均误差刚好为 5\%，说明误差约束是决定最优解的重要边界。
\end{enumerate}

最终，第二问方案在满足总处理量不低于 1、加权平均误差不超过 5\%
的条件下，将日电费降低到 162.3362 元，相比问题一固定方案节约 28.0318
元，节约率为 14.73\%。

\begin{center}\rule{0.5\linewidth}{0.5pt}\end{center}

\subsection{参考文献}\label{ux53c2ux8003ux6587ux732e}

\subsection{一、实时系统误差容忍和约束相关}\label{ux4e00ux5b9eux65f6ux7cfbux7edfux8befux5deeux5bb9ux5fcdux548cux7ea6ux675fux76f8ux5173}

\subsubsection{\texorpdfstring{1. \emph{Real-time optimization of
average error under
uncertainty}}{1. Real-time optimization of average error under uncertainty}}\label{1-real-time-optimization-of-average-error-under-uncertainty}

\begin{itemize}
\item
  Bernstein, A., Bouman, N. J., \& Le Boudec, J.-Y. (2016).

  Real‑Time Minimization of Average Error in the Presence of Uncertainty
  and Convexification of Feasible Sets

  .

  arXiv

  .

  \begin{itemize}
  \item
    该文讨论实时控制系统中误差处理的方法，考虑在实时约束条件下最小化误差的平均值，并给出了误差动态约束下的可行集构造。可用于支持：

    \begin{itemize}
    \item
      实时误差控制 vs. 平均误差优化的区分
    \item
      在实时控制问题中误差约束的重要性
    \end{itemize}
  \item
    （可引用于论文建模误差约束讨论部分）
  \end{itemize}
\end{itemize}

\subsubsection{\texorpdfstring{2. \emph{Average Task Execution Time
Minimization under
(m,k)‑constraints}}{2. Average Task Execution Time Minimization under (m,k)‑constraints}}\label{2-average-task-execution-time-minimization-under-mkconstraints}

\begin{itemize}
\item
  Shi, J. (2023).

  Average Task Execution Time Minimization under (m, k)-soft error
  constraints

  (RTAS Workshop).

  \begin{itemize}
  \item
    这类研究提出了在 \textbf{(m, k)-误差限容约束}
    情况下调度策略，其中允许\textbf{某些任务出现错误但保证大部分正确}，这正是实时任务对``部分误差可容忍''的数学规则。
  \item
    例如要求在 k 个时段内至少有 m 个误差低于阈值，这与瞬时 vs.
    平均误差之间的权衡一致，可作为建模假设依据。
  \end{itemize}
\end{itemize}

\subsubsection{\texorpdfstring{3. \emph{Real-Time Error Control for
Simulation}}{3. Real-Time Error Control for Simulation}}\label{3-real-time-error-control-for-simulation}

\begin{itemize}
\item
  Bui, H. P., Tomar, S., Courtecuisse, H., \& Bordas, S. (2016).

  Real-Time Error Control for Surgical Simulation

  .

  arXiv

  .

  \begin{itemize}
  \item
    该文展示了在实时模拟系统中如何控制误差不超过可容忍水平，是实时系统误差约束的一种实例参考。
  \item
    支持``实时控制系统误差要求通常是
    \emph{局部/每步不超过阈值}''的说法。
  \end{itemize}
\end{itemize}

\begin{center}\rule{0.5\linewidth}{0.5pt}\end{center}

\subsection{二、批处理 /
平均误差评估相关}\label{ux4e8cux6279ux5904ux7406--ux5e73ux5747ux8befux5deeux8bc4ux4f30ux76f8ux5173}

虽然不完全同你题目的一对一对应，以下论文可以作为``对于批处理/平均误差衡量普遍采用
\emph{平均} 或 \emph{累积} 误差评估''的侧面支持：

\subsubsection{\texorpdfstring{4. \emph{Performance Metrics in Machine
Learning}}{4. Performance Metrics in Machine Learning}}\label{4-performance-metrics-in-machine-learning}

\begin{itemize}
\item
  Botchkarev, A. (2018).

  Performance Metrics (Error Measures) in Machine Learning Regression,
  Forecasting and Prognostics: Properties and Typology

  .

  arXiv

  .

  \begin{itemize}
  \item
    这篇综述总结了误差评估指标（包括平均误差、加权平均误差等），说明：

    \begin{itemize}
    \item
      在机器学习或批处理分析任务中，更常用\textbf{平均误差或整体性能指标}评估模型，而不是严格每次输出误差。
    \end{itemize}
  \item
    可用于支持``对于批处理/数据分析任务通常采用平均或累积误差约束''的说法。
  \end{itemize}
\end{itemize}

\subsubsection{\texorpdfstring{5. \emph{Challenges of Managing Errors in
Batch Processing
Systems}}{5. Challenges of Managing Errors in Batch Processing Systems}}\label{5-challenges-of-managing-errors-in-batch-processing-systems}

\begin{itemize}
\item
  Jangam, S. K., \& Pedda Muntala, P. S. (2023).

  Challenges and Solutions for Managing Errors in Distributed Batch
  Processing Systems and Data Pipelines

  .

  Int. Journal of Emerging Research in Engineering and Technology

  .

  \begin{itemize}
  \item
    该综述讨论了批处理系统错误管理、容错和误差质量保证问题，说明分布式批处理系统通常关注\textbf{总体数据质量
    / 平均误差控制}，而不是每个时间步满足误差。
  \end{itemize}
\end{itemize}

\begin{center}\rule{0.5\linewidth}{0.5pt}\end{center}

\subsection{三、实时系统反馈与控制背景}\label{ux4e09ux5b9eux65f6ux7cfbux7edfux53cdux9988ux4e0eux63a7ux5236ux80ccux666f}

这些论文与实时控制系统误差性质有关，可用于补充\textbf{实时误差绝对约束}的背景说明：

\subsubsection{\texorpdfstring{6. \emph{Minimum Entropy Control of
Tracking
Errors}}{6. Minimum Entropy Control of Tracking Errors}}\label{6-minimum-entropy-control-of-tracking-errors}

\begin{itemize}
\item
  Yue, H., \& Wang, H. (2003).

  Minimum Entropy Control of Closed‑Loop Tracking Errors for Dynamic
  Stochastic Systems

  ,

  IEEE Trans. on Automatic Control

  .

  \begin{itemize}
  \item
    讨论在实时控制系统中控制跟踪误差的数学方法，是实时系统误差控制的经典参考。
  \end{itemize}
\end{itemize}

\subsubsection{\texorpdfstring{7. \emph{Real-Time Fault
Tolerance}}{7. Real-Time Fault Tolerance}}\label{7-real-time-fault-tolerance}

\begin{itemize}
\item
  Persya, A. C., \& Nair, T. R. G. (2010).

  Fault Tolerance in Real Time Multiprocessors -- Embedded Systems

  .

  arXiv

  .

  \begin{itemize}
  \item
    研究实时多处理器系统中任务执行误差与容错问题，强调实时系统对错误容忍性与调度的严格度。
  \end{itemize}
\end{itemize}

\begin{center}\rule{0.5\linewidth}{0.5pt}\end{center}

◆ \textbf{实时实时系统误差控制习惯}：

\begin{itemize}
\item
  Yue \& Wang (2003) -- 实时误差控制的数学理论
\item
  Bui et al. (2016) -- 实时仿真误差控制案例
\item
  Persya \& Nair (2010) -- 实时多核/嵌入式误差容忍与调度
\end{itemize}

◆ \textbf{批处理/数据分析任务误差约束方式}：

\begin{itemize}
\item
  Botchkarev (2018) -- 平均误差 / 综合误差指标在数据分析中的应用
\item
  Jangam \& Pedda Muntala (2023) -- 批处理系统误差管理综述
\end{itemize}

◆ \textbf{软实时 / 容错调度约束建模技巧}：

\begin{itemize}
\item
  Shi (2023) -- (m, k)-约束模式，介于严格实时和平均误差间
\end{itemize}

\end{document}
